\documentclass[../main.tex]{subfiles}

\begin{document}
% -----------------------------
\section{Dead code elimination}
% -----------------------------
Next, will add a tiny DCE pass that removes ops whose results have no use.
Focus here is to erase ops that:
\begin{itemize}
    \item are not terminators (e.g., not func.return)
    \item have results
    \item all results have zero uses
    \item op is in a “safe-to-drop” set 
\end{itemize}
We will run this in a loop until no more ops can be removed.
In most xDSL versions, the safe sequence is:
\begin{itemize}
    \item op.detach() (removes it from its block / unlinks)
    \item then op.erase() (frees it, asserts parent is None)
\end{itemize}

First, we will create some helper functions. Start with the function that returns how many times an SSA value is used in the IR.
\begin{lstlisting} [language=Python, caption={Counting uses}]
# -------------------------------------
# Dead code elimination
# -------------------------------------
def _num_uses(v: SSAValue) -> int:
    uses = getattr(v, "uses", None)
    if uses is None:
        return 0
    try:
        return len(uses)
    except TypeError:
        # Sometimes uses is an iterator/generator-like
        return sum(1 for _ in uses)
\end{lstlisting}
In xDSL (like in MLIR), every SSAValue keeps track of its uses — i.e. which operations consume that value as an operand.
Next, check conditions for the simple dce case:
\begin{lstlisting} [language=Python, caption={Checking if it is trivial to do dce}]
def _is_trivially_dceable(op: Operation) -> bool:
    """
    Conservative: only remove arith ops (constants and arithmetic) that have
    no users. Not covering memref/loads/stores/calls etc. for now.
    """
    if op.name == "func.return":
        return False

    # Only drop ops from these dialects for now
    if op.name.startswith("arith."):
        return True

    # If we also want to also drop leftover hc.* after lowering, we could include:
    # if op.name.startswith("hc."):
    #     return True

    return False
\end{lstlisting}
Implement a function that performs one pass of Dead Code Elimination (DCE) over a single Block in xDSL:
\begin{lstlisting} [language=Python, caption={One DCE pass implementation}]
def dce_block(block: Block) -> bool:
    """
    One pass of DCE over a block.
    Returns True if anything was erased.
    """
    changed = False

    # Iterate backwards: safer when erasing
    for op in list(block.ops)[::-1]:
        # Must be safe to erase
        if not _is_trivially_dceable(op):
            continue

        # If op has no results, keep it (could be side-effecting; conservative)
        if len(op.results) == 0:
            continue

        # If any result is used, keep it
        if any(_num_uses(res) > 0 for res in op.results):
            continue

        # Otherwise erase
        _erase_op(op)
        changed = True

    return changed
\end{lstlisting}
Notice that by iterating backwards, we eliminate leaf nodes first.
\\Implement a function that runs Dead Code Elimination (DCE) repeatedly over all
function bodies in a module until nothing more can be removed.
\begin{lstlisting} [language=Python, caption={Applying DCE}]
def apply_dce(module) -> None:
    """
    Repeatedly run DCE on all function body blocks until fixpoint.
    """
    changed = True
    while changed:
        changed = False
        for top_block in module.body.blocks:
            for op in list(top_block.ops):
                if isinstance(op, func.FuncOp):
                    for body_block in op.body.blocks:
                        if dce_block(body_block):
                            changed = True
\end{lstlisting}
Let’s break it down:
\begin{itemize}
    \item It walks through all functions in the module
    \item Runs \path{dce_block()} on every block inside each function
    \item If anything was removed → run everything again
    \item Stops only when no more dead code can be eliminated
\end{itemize}
Lastly, implement a helper function that first removes operation from its parent block,
and then erases it.
\begin{lstlisting} [language=Python, caption={Erasing operation}]
def _erase_op(op: Operation) -> None:
    # Detach from parent block first (required in our xdsl)
    if hasattr(op, "detach"):
        op.detach()
    else:
        raise RuntimeError("Operation.detach() not found; cannot safely DCE ops.")
    op.erase()
\end{lstlisting}
Now, use it in the main function:
\begin{lstlisting} [language=Python, caption={Adding DCE in the main body}]
    ...
    apply_dce(module)
    print("\n=== AFTER DCE ===")
    print(module)
\end{lstlisting}
In the DCE section, you should only see one constants, others have been removed:
\begin{lstlisting} [language=bash, caption={DCE output}]
=== AFTER DCE ===
builtin.module {
  func.func @my_func() -> i32 {
    %0 = arith.constant 0 : i32
    func.return %0 : i32
  }
}
\end{lstlisting}

\end{document}
